
\newcommand{\rChist}[5]{\item[#1] \emph{Old:} #2\par\emph{New:} #3\par\emph{Reason:} #4\par\emph{Files:} {\small #5}}

\section{Superseded estimates, failed approaches and retracted statements}\label{rep:history}

\begin{center}
\fbox{\parbox{0.9\textwidth}{\centering\textbf{Superseded / not used in the papers.} Every statement and number in the ``Old'' fields of this part is superseded and must not be quoted as a result. The current results are those of Parts~I and~II and of the Letter and the paper.}}
\end{center}

This part records how the results changed during the research: earlier estimates that later calculations replaced, approaches that failed or were abandoned, and statements that were retracted.
It serves two purposes.
A reader who meets an older number in a talk, a draft or the archived research files can find what replaced it and why; and the record shows which conclusions changed under scrutiny and which survived it.
Each item gives the old statement, the new one, the reason, and the files of the research tree that document the change (derivation files D$nnn$, calculations C$nnn$, review reports R$nn$ and the blind checks of Sec.~\ref{rep:v-blind}).
``The internal draft'' is the version of the paper that preceded the present papers; it was not released.
The ``New'' fields quote current numbers only to identify the replacement; their full statements and conditions are in the sections cited.

\subsection{The Higgs zero mode}\label{sec:v23-higgs}

\begin{description}
\rChist{H1}{The internal draft assumed that new physics keeps the effective quartic coupling positive at $1/L_b$ and took $\lambda=0.01$, a ``stabilized Higgs''.}
{The Standard Model itself has $\lambda(\mub)\approx-0.0098<0$; the Higgs zero mode is a finite-volume resonance confined by the zero-point energy of the gauge zero modes down to the threshold $\lambda_c\approx-0.012$ ($\Pone=1/2$), and the Standard Model lies on the safe side for the measured top mass (Secs.~\ref{rep:h-matching}--\ref{rep:h-smmap}).}
{The second stage showed that the constant gauge modes confine the zero mode at negative $\lambda$; no new physics is needed below $\mub$.}
{internal draft (main\_v11.tex l.171); V12\_RESULTS\_INDEX \S E (E1$'$, E3$'$); D110, D112, D121, D122; run05.}

\rChist{H2}{``$\lambda\simeq-0.014$ near the Planck scale'' as the relevant coupling.}
{$\lambda(\mub)=-0.0098$ at $\mub=\MPbar/722$; the Planck-scale value is not the relevant one.}
{The matching scale is $1/L_b$, not the Planck mass.}
{D120.}

\rChist{H3}{``A non-minimal coupling does not act at the turn'', because $RL_b^2=12(hL_b)^2\approx2\times10^{-4}$.}
{True at $\tau=0$, but the late part of the turn ($\tau\approx3.5$--$6$) is decided by $\xiH$: with $\xiH=0$ the zero mode runs away, and survival requires $\xiH(\mub)\ge\xi_c(\Pone)$, about $0.074$--$0.08$ at the measured mass (Sec.~\ref{rep:h-universality}).}
{The independent check R14 found that the term $a''/a$ makes the width grow from $1.9\times10^{-3}$ at $\tau=3$ to $1.15$ at $\tau=5$ for $\xiH=0$.}
{internal draft l.171; R14 correction addendum to D108 \S2; D111; D129.}

\rChist{H4}{Inputs of the internal draft: $\mu=1/L_b=3.10\times10^{15}$~GeV; one-loop couplings $g_2/g_3/g_Y=0.5268/\allowbreak 0.5396/\allowbreak 0.4415$ and $y_t=0.4585$; $\alpha_s(M_Z)=0.1184$.}
{$\mub=3.3726\times10^{15}$~GeV; NNLO couplings $0.5278/\allowbreak 0.5357/\allowbreak 0.4428$ and $y_t(\mub)=0.4278$; $\alpha_s=0.1180$ (Sec.~\ref{rep:h-smmap}). The one-loop couplings survive only in the DMC shape correction, where their effect is the $\pm0.002$~GeV item of the budget.}
{Updated inputs (PDG 2026) and the NNLO mapping with SMDR-type matching.}
{internal draft l.546, l.558; sm\_provider.py; D120 \S3.6.}

\rChist{H5}{Zero-mode mass $m_\Phi^2=(\frac94g_2^2+\frac34g_Y^2+6\lambda)T_B-6y_t^2T_F$ with ``three polarizations per massive vector''.}
{Consistent one-loop matching in Coulomb gauge with $\MSbar$ finite-volume sums: two transverse and one Coulomb mode per massive vector, mixed with the Goldstone; exact off-shell Goldstone mass (Sec.~\ref{rep:h-matching}).}
{The unitary-gauge counting is not a consistent Born--Oppenheimer split; it made the result depend on the gauge by about $0.25$~GeV (R14).}
{internal draft l.556; D110; R15 check of D110.}

\rChist{H6}{First stage: main model ``MU''; critical masses $172.22/\allowbreak 172.59$~GeV (D108) or $172.20/\allowbreak 172.57$~GeV (R14); $\Pone(172.60)\approx0.48$ ($0.36$--$0.74$); ``$\xiH\gtrsim1/6$ and $m_t\lesssim172.2$~GeV give $\Pone\ge0.9$''; scheme error $\sim0.25$~GeV; reduction error $\pm0.1$--$0.15$~GeV; DMC bracket $+0.04$ to $+0.23$~GeV.}
{All replaced by the second and third stages; the final critical masses are $172.502/\allowbreak 172.960/\allowbreak 173.418$~GeV (Sec.~\ref{rep:h-survival}).}
{Inconsistent matching, one-dimensional reduction and no full-dimensional validation.}
{D108; V12\_RESULTS\_INDEX \S A (note).}

\rChist{H7}{``Freezing the width beyond $\tau=3$ lowers $P$ by only about $0.01$''.}
{For $\xiH=0$ the width rises from $1.9\times10^{-3}$ to $1.15$ between $\tau=3$ and $5$; all later hazards use $\xiH=1/6$ on $\tau\in[0,5]$ with the real-time condition on $\xiH$.}
{Independent check R14.}
{R14 correction addendum to D108 \S2.}

\rChist{H8}{First-stage thresholds on $\xiH$: $\xi_c\approx0.16$ (resonance measure) or $0.04$--$0.12$ (classical runaway measure).}
{A single definition, $P_{\rm surv}(\xi_c)=0.99\,P_{\rm surv}(1/6)$; $\xi_c$ is a universal function of $\Pone$ (Sec.~\ref{rep:h-universality}).}
{Two measures gave inconsistent answers; real-time evolution decides.}
{D111 \S0(3); D129.}

\rChist{H9}{Second stage (D110): consistent matching gives critical masses $172.35/\allowbreak 172.77$~GeV and $\Pone(172.6)=0.70$; ``the zero-mode cross terms stabilize by $+0.04$ to $+0.08$~GeV (not included)'', widened after the R15 check to $+0.04$ to $+0.09$~GeV.}
{The cross terms destabilize: $-0.080/\allowbreak{-}0.118/\allowbreak{-}0.160$~GeV at $\Pone=0.9/\allowbreak 0.5/\allowbreak 0.1$, included in the model (Sec.~\ref{rep:h-cross}).}
{The estimate used Higgs-only, periodic-boson curvatures; with all fields and antiperiodic fermions the hypercharge curvature of the light fermions lifts the well more than the barrier. Confirmed blindly by V2.}
{D110; D121; R15 check of D112; R16 V2.}

\rChist{H10}{D112: critical masses $172.51/\allowbreak 173.00$~GeV from width ratios; ``the soft hypercharge direction is unstable for $r<0.776$ and needs periodic holonomies''.}
{Direct two-dimensional roots at every node give $172.518/\allowbreak 173.014$~GeV (same model); the instability is an artifact of keeping only Higgs fluctuations.}
{Node-direct calculation (no ratio transfer); full set of fields in the curvature.}
{D122 \S4.3; D121 \S0.2; R16 V2.}

\rChist{H11}{Second-stage final table (D116): $172.517/\allowbreak 173.012/\allowbreak 173.514$~GeV; $\Pone(172.60)=0.856$; $\lambda_c=-0.00965/\allowbreak{-}0.01291/\allowbreak{-}0.01624$; $m_c-m_0=1.49/\allowbreak 1.99/\allowbreak 2.49$~GeV; $\xi_c=0.036/\allowbreak 0.072/\allowbreak 0.120$; inflation-scale item $\pm0.05$~GeV; Standard Model mapping error $\pm0.16$--$0.17$~GeV; cross terms ``$+0.04/\allowbreak{+}0.06$~GeV, one-sided''.}
{$172.502/\allowbreak 172.960/\allowbreak 173.418$~GeV; $\Pone(172.60)=0.843$; $\lambda_c=-0.0091/\allowbreak{-}0.0121/\allowbreak{-}0.0152$; $m_c-m_0=1.41/\allowbreak 1.87/\allowbreak 2.33$~GeV; $\xi_c=0.035$ ($172.2$~GeV), $0.074$ ($172.6$~GeV), $0.12$ ($\Pone=1/2$); inflation-scale item $\pm0.035/\pm0.067/\pm0.112$~GeV; Standard Model error $\pm0.07$~GeV ($\mtbar$) and $\pm0.10$~GeV (pole); cross terms $-0.08/\allowbreak{-}0.12/\allowbreak{-}0.16$~GeV (Secs.~\ref{rep:h-survival}--\ref{rep:h-universality}).}
{NNLO mapping, node-direct two-dimensional hazard, cross terms, blind error analysis of the mapping (V1).}
{D116; D124 \S4; V12\_RESULTS\_INDEX \S E; numbers\_run05.csv (R5-001, R5-029, R5-030, R5-042--R5-047).}

\rChist{H12}{D111: $\xi_c=0.032/\allowbreak 0.054/\allowbreak 0.120$ (MU potential); ``$\xiH\gtrsim0.03$--$0.12$, of the same order as the inflation bound $\xi\gtrsim0.06$''; stochastic nucleation with $B=192\pi^2\xi^2/|\lambda|$ ``matters only for $\xi\lesssim0.020$--$0.028$''.}
{Current $\xi_c$ of Sec.~\ref{rep:h-universality}; compared at the same scale, the zero-mode condition is stronger than the inflation bound for $m_t\gtrsim172.1$~GeV; the Hawking--Moss coefficient is $96\pi^2$, which gives $0.028$--$0.040$ for the MU potential.}
{Comparison at mixed scales; a factor of two in the Hawking--Moss exponent.}
{D129 \S0.4, \S5; D127 \S7; R16 V5, V6.}

\rChist{H13}{D111: a real Gaussian initial state gives $w_0=0.53$--$0.93$.}
{That state, truncated at the origin, is not in the domain of the Hamiltonian; the time-symmetric minimal-variance rule gives $w_0=0.98$--$0.995$ (Sec.~\ref{rep:h-survival}).}
{Independent check R15 of D111.}
{D124 \S3; R15 addendum to D111 \S2.}

\rChist{H14}{Standard Model mapping errors: run01 $\pm0.13$~GeV (pole) and $\pm0.11$~GeV ($\mtbar$); run02 $\pm0.086$ and $\pm0.047$~GeV.}
{$\pm0.10$~GeV (pole) and $\pm0.07$~GeV ($\mtbar$), four-loop truncated pole relation (Sec.~\ref{rep:h-smmap}).}
{The blind check V1 explained the factor-of-ten difference of matching-scale variations between SMDR and \texttt{mr} and recommended these values.}
{D120; R16 V1 (DIAGNOSIS.json \texttt{recommended}).}

\rChist{H15}{D120 run01, without cross terms: $172.580/\allowbreak 173.074/\allowbreak 173.575$~GeV; $\mtbar=162.80/\allowbreak 163.27/\allowbreak 163.75$~GeV.}
{$172.502/\allowbreak 172.960/\allowbreak 173.418$~GeV; $\mtbar=162.73/\allowbreak 163.16/\allowbreak 163.60$~GeV.}
{Cross terms and their one-loop completion.}
{R3\_mbar\_run05.json; R16 consistency audit.}

\rChist{H16}{Cross terms by the ratio method: $-0.093/\allowbreak{-}0.135/\allowbreak{-}0.180$~GeV; a first vev-mixing run with a coefficient four times too large; ``the top term is at most $25\%$ of the vev mixing''.}
{Direct roots: $-0.085/\allowbreak{-}0.124/\allowbreak{-}0.167$~GeV (run02), and $-0.080/\allowbreak{-}0.118/\allowbreak{-}0.160$~GeV after the covariant vertex, the Coulomb mismatch and the exact top chiral loop (run05); the vev-mixing and top-chiral groups are comparable and largely cancel.}
{Ratio transfer replaced by direct roots; error in the first vev-mixing vertex; exact top bubble.}
{D121 run02--run05 (\texttt{*\_factor4error\_superseded.json}).}

\rChist{H17}{First claims layer (D126 v1): $\Pbar\approx0.79$ ($[0.53,0.995]$); FCC-ee exclusion at $173.8$~GeV; $M_h=125.20$~GeV; ``implied by MMR 2022''.}
{$\Pbar=\PbarOne$ (direct mass); exclusion reach $\mExclFCC$~GeV; $M_h=125.13$~GeV; the bound of Mantziris, Markkanen and Rajantie (2021) applies (Secs.~\ref{rep:h-budget}, \ref{rep:h-universality}).}
{Wrong Higgs mass, wrong reference, earlier centre.}
{D126 v2; R16 V6.}

\rChist{H18}{D126 v2 before correction: variant of Bednyakov et al.\ $0.626$; $\MSbar$ cross-section value $0.557$; $B_{\rm sel}=1.459$ for a prior stated as flat on $[150,200]$~GeV; ``$|\ln B|<0.8$''; HL-LHC separation $\approx1.6\sigma$.}
{Equal-mass convention throughout ($0.67$; quoted range $\PbarBKPV$); $\PbarMbarSigma$; $B_{\rm sel}=\BselDirect$ with upper bound $\BselMax$; $|\ln B|\le\BayesMax$; HL-LHC $\sigHLLHCmin$--$\sigHLLHCmax\sigma$ (final values, Sec.~\ref{rep:h-budget}).}
{Blind check V3: two conventions for asymmetric errors in one code; the prior was applied on a grid over $[160,186]$~GeV; the HL-LHC reach depends on the assumed precision of $\alpha_s$.}
{D126 v2 correction addendum; R16 V3.}

\rChist{H19}{Literature notes: ``MMR 2022 ($\xi\gtrsim0.1$) is stronger than $\xi_c$ for $m_t\le172.6$~GeV''; a reading in which ``the inflation bound is stronger below $172.4$~GeV''.}
{The plateau is a minimally coupled scalar, so MMR 2021 applies ($\xi_{\rm EW}\approx0.049$--$0.060$); the second reading compared different scales; at the same scale the zero-mode condition is stronger for $m_t\gtrsim172.1$~GeV.}
{Blind literature check V6 and blind recomputation V5.}
{D127 \S0.1; D129 \S0.5; R16 V5, V6.}

\rChist{H20}{``$-0.19$~GeV per e-fold of the duration of the turn''.}
{Two correct measures with different definitions: the secant over one e-fold, $-0.170/\allowbreak{-}0.270/\allowbreak{-}0.394$~GeV, and the local derivative, $-0.19/\allowbreak{-}0.31/\allowbreak{-}0.47$~GeV (Sec.~\ref{rep:h-budget}).}
{Blind check V3.}
{D124; D128; R16 V3.}

\rChist{H21}{Second stage: ``the two-dimensional to one-dimensional width ratio is $1.06$--$1.08$ and needs a bump $A=0.0937$ outside the well''.}
{In the final model the ratio is $1.010$--$1.018$ and the bump $0.013$.}
{Model refinements (matching, cross terms) removed most of the difference.}
{D129 \S0.1.}

\rChist{H22}{Earlier claim wording: ``$\Pbar\approx0.8$'', ``$\xiH\gtrsim0.07$'', ``exclusion needs $m_t\gtrsim173.9$~GeV''.}
{$\Pbar\approx\PbarOne$; $\xiH(\mub)\gtrsim\xiHmin$ (conservative); $95\%$ exclusion of the theory curve at $\mExclTheory$~GeV and FCC-ee reach $\mExclFCC$~GeV.}
{Final model and corrected claims layer.}
{D126 v2 \S1.4, \S4; D132.}

\rChist{H23}{Error-budget items of the third stage: reduced model ``in-model'' $^{+0.009}_{-0.012}$, $^{+0.012}_{-0.017}$ and $^{+0.015}_{-0.022}$~GeV with DMC $\pm1\sigma$ and $B_x$ sub-items; full-dimensional item $+0.025$~GeV at $\Pone=0.1$ by extrapolation; matching item from the one-dimensional scale scan of the second stage ($+0.014/\allowbreak{-}0.045$~GeV at $\Pone=1/2$); preparation item $-0.016$~GeV at all levels.}
{Reduced-model item $\pm0.002$~GeV (fixed couplings in $\delta E_{\rm DMC}$); full-dimensional item $+0.012/\allowbreak{+}0.020/\allowbreak{+}0.031$~GeV computed with the final slopes (Sec.~\ref{rep:v-fulldim}); matching item from the scan in the final model, $^{+0.009}_{-0.039}$, $^{+0.012}_{-0.049}$ and $^{+0.015}_{-0.058}$~GeV; preparation $-0.016/\allowbreak{-}0.006/\allowbreak{-}0.004$~GeV (Sec.~\ref{rep:h-budget}).}
{The full-dimensional check measured directly what the bracketed sub-items estimated; overlap removed; all items recomputed in the final model.}
{table1\_rebuild.json; numbers\_run05.csv R5-035--R5-039 (column \texttt{supersedes}); D123; D132.}

\rChist{H24}{First-stage renormalization-group improvement of the zero-mode potential with $\mu^2=(1+r^2)/L^2$ (about $+0.13$~GeV).}
{Removed: the finite-volume one-loop energy contains no $\ln(m^2L^2)$, and the masses are far below $2\pi/L$ (Sec.~\ref{rep:h-matching}).}
{The infinite-volume logarithm cancels against the finite-volume sum in the consistent matching.}
{D110 \S3.5.}

\rChist{H25}{The one-loop vertex of the cross terms in ordinary Coulomb gauge (anticommutator vertex) with the top chiral loop: $172.495/\allowbreak 172.950/\allowbreak 173.407$~GeV; the gauge logarithms closed only to $61\%$, and a ``missing $39\%$'' item of $\pm0.009/\pm0.012/\pm0.014$~GeV gave a combined cross-term error of $\pm0.012/\pm0.016/\pm0.019$~GeV.}
{The background-covariant Coulomb-gauge vertex with the Coulomb mismatch closes the logarithms to $97$--$99\%$; $172.502/\allowbreak 172.960/\allowbreak 173.418$~GeV with a model-internal error of $\pm0.008/\pm0.010/\pm0.013$~GeV (Sec.~\ref{rep:h-cross}). The convergence table keeps the intermediate row only to show the size of the step.}
{The vertex must be in the same gauge as the twisted transverse modes of the conditional constant approximation.}
{D121 appendices run04, run05.}

\rChist{H26}{First vev-mixing run of the cross terms: diagonal channels only, a polarization factor between $2/3$ and $1$, ``$+0.010$ to $+0.026$~GeV, stabilizing''.}
{Four channels with the exact factor $2/3$ and a top chiral loop of comparable size (run04), then the covariant vertex (run05); see H16 and H25.}
{Omitted channels and the choice of gauge.}
{D121 appendices run03, run04.}

\rChist{H27}{Second-stage Standard Model mapping: two-loop running with a constant three-loop calibration and fitted matching; Standard Model error ``$\pm0.16$--$0.17$~GeV from the literature''.}
{NNLO matching of SMDR type with three-loop running and four- and five-loop QCD terms; it moves all critical masses by $+0.063$~GeV (Sec.~\ref{rep:h-smmap}).}
{The fitted matching gave too large a $g_3$ and lacked the four-loop QCD terms of $y_t$.}
{D120 \S3.6.}

\rChist{H28}{Matching-scale error of the mapping from a renormalization-group reconstruction of the three-loop logarithms, $[-0.088,+0.102]$~GeV.}
{The compiled SMDR with independent scale variations gives $[-0.042,+0.054]$~GeV without the pure-QCD tower; the reconstruction overestimates by a factor of about three and is kept only as a comparison (Sec.~\ref{rep:h-smmap}).}
{Direct computation with the compiled SMDR became available.}
{D120 \S4, App.~R.4.}

\rChist{H29}{Draft wording: ``the absolute stability bound is $171.5$~GeV and $m_0$ lies below it''.}
{With the same inputs $m_{\rm stab}=170.99$~GeV lies below $m_0=171.09$~GeV (Sec.~\ref{rep:h-smmap}).}
{The order was wrong; checked with the same mapping.}
{package \texttt{higgs/01\_sm\_mapping/abs\_stability.json}; consistency check of the draft against this file (2026-10-16).}

\rChist{H30}{``$P_2=\Pone^2$ without condition'' for the survival of both branches.}
{The two branches interfere maximally; there is no decohering history probability, only Fr\'echet bounds $\max(0,2\Pone-1)\le P_2\le\Pone$, and the linear value $\Pone^2/w_0$ is one choice (Sec.~\ref{rep:h-survival}).}
{Gell-Mann--Hartle analysis of the two-time histories.}
{D108; D130.}

\rChist{H31}{Draft table: ``$hL_b\times1.1$: $-0.020/\allowbreak -0.032/\allowbreak -0.049$~GeV''; draft prefactor ``$A\approx300$'' of the closed-form threshold.}
{The shifts $-0.020/\allowbreak{-0.032}/\allowbreak{-0.049}$~GeV belong to $hL_b\times0.9$, and $\times1.1$ gives $+0.018/\allowbreak{+0.030}/\allowbreak{+0.046}$~GeV; the effective prefactor is $A_{\rm eff}\simeq310$ (Sec.~\ref{rep:h-reduced}).}
{Sign error and an inconsistent prefactor found when the papers were checked against the result files.}
{R16 V3 \texttt{sens\_results.json}; D128 \S1.2; consistency check of the drafts against the result files (2026-10-16).}
\rChist{H32}{Zero-mode energy of the earlier draft computed with a decoupled quartic part, the ``stabilized'' $\lambda=0.01$ and $y_t=0.4585$.}
{The Higgs zero mode is a finite-volume resonance at the Standard Model coupling $\lambda(\mub)=-0.0098$ (Secs.~\ref{rep:h-matching}--\ref{rep:h-survival}); the zero-mode share evaluated with the real part of the resonance energy is $0.298$, and the baseline $\fzm=0.30$ is kept (Sec.~\ref{rep:zeromodes}).}
{Consequence of H1 for the zero-mode energy; the change of $\fzm$ is below $1\%$.}
{D108 \S5; D124; long PRD draft appendices on the Casimir energy and the zero modes.}\end{description}

\subsection{The emptiness gate, the shape and the time symmetry}\label{sec:v23-shape}

\begin{description}
\rChist{S1}{``The rest is worth fewer than five quanta''; ``a single quantum with $|n|\ge5$ already forbids the turn''.}
{About five quanta: $\sum|n|<4.76$ lowest bosonic quanta ($4.70$--$5.05$ over the range of $\fzm$); the single-quantum statement holds only for $\fzm>0.265$, and the statement that holds for all $\fzm>0.250$ is $|n|\ge\sqrt{26}$ (Sec.~\ref{rep:gate}).}
{Recomputed budget with the zero-mode energy; the bound equals five at $\fzm=0.265$.}
{internal draft l.28, 42, 224, 418; D125 \S7; R16 V4.}

\rChist{S2}{A real Gaussian shape state prepared at the turn as a premise; ``the state must be time symmetric''; ``the two branches are mutually exclusive''.}
{A theorem: any shape state is real at most once, and, for a real Gaussian state, survival of the shape on both branches with $\ge90\%$ ($99\%$) probability requires that moment within $0.35/h$ ($0.24/h$) of the turn; the equivalences B1/B2 hold under stated conditions (Sec.~\ref{rep:quantum-shape}).}
{The premise was replaced by a proved statement and a conditional theorem; the time symmetry is a conclusion under conditions, not a requirement.}
{internal draft l.28, 42, 44, 324; D109; D113; D125; D131.}

\rChist{S3}{Survival thresholds: $s_{99}=2.2$, $\sigma(99\%)\approx0.24$; ``survival $\ge99/90/50\%$ requires $|t_x|\lesssim0.2/\allowbreak 0.38/\allowbreak 0.62/h$'' (interpolated); first runs $t_*(99/90\%)=0.232/\allowbreak 0.340$ ($N=256$--$512$); next order $s_{99}=2.02$, $\sigma(99\%)\approx0.20$.}
{$s_{99}\approx2.8$--$2.95$; linear $99\%$ threshold $\approx0.29$; $t_*=0.237/\allowbreak 0.355/\allowbreak 0.631$ (direct, $N=2048$); next order $s_{99}=3.16$, $\sigma\approx0.30$ (Sec.~\ref{rep:quantum-shape}).}
{Too few samples for $99\%$-level quantities; the blind checks T19 and V4 found the larger values, and the main calculation, repeated with $N=2048$, confirmed them.}
{D114; D125 run01, run02; R16 T19, V4.}

\rChist{S4}{``A real state set in the past frozen region gives $\sigma_E\gtrsim2.0$''; ``the shape must be prepared within about $0.1/h$ of the turn, $70$ times more strictly than the light fields''.}
{Such states destroy the shape in $82$--$87\%$ of cases and the survivors have $\sigma_E\approx0.7$--$0.9$; the latest preparation is $0.24$--$0.36/h$, $21$--$32$ times stricter in Hubble times ($65$--$135$ times in volume e-folds).}
{Nonlinear survival computed; the R14 check corrected the factor in the freezing region.}
{D109 \S7; D114; D125 \S3.2; R14 check of D109.}

\rChist{S5}{The criterion of Proposition~2 ``$<1$'', ``violated for $N_{\rm pre}\gtrsim1.5$--$2$''.}
{The criterion is $<3^{-3/2}=0.19245$; it is violated at $1.32$ ($c=1$) and $1.01$ ($c=3$).}
{Corrected derivation of the shear criterion.}
{D114 \S5; D125 \S6.}

\rChist{S6}{B1 stated without condition; B3 ``minimal energy variance'' without restriction; ``a cap is impossible in this model''.}
{B1 requires a single classical history; the Wheeler--DeWitt test makes this concrete: the Airy components $\mathrm{Ai}$ and $\mathrm{Bi}$ must carry the same shape state, and decay under the barrier is sufficient (if they carry different real states, the reality defect tends to $2\eta^2/(1+\eta^2)$); minimal variance holds within the centred Gaussian family; the divergence of the homogeneous stress at a cap is a limit of the model's precision, and the curved-cigar sign theorem shows that curvature does not flip the aspect kernel.}
{Independent check R15 of D113 (counterexample $\mathrm{Ai}(-\zeta)\varphi_0+\mathrm{Bi}(-\zeta)\varphi_1$).}
{D113 and its addendum; D118; D131; R15 check of D113.}

\rChist{S7}{``Every draw survives up to three times the natural amplitude''.}
{True only for free fields; with the zero-mode energy, three times the amplitude survives $97.3\%$ ($N=256$), $98.3\%$ ($N=2048$) and $98.7\%$ (T19).}
{The free-field statement was carried over to the Standard Model with zero modes.}
{internal draft l.362, l.617; D114; D125 run02; R16 T19.}

\rChist{S8}{``The Moyal correction is expected at the percent level; we have not computed it.''}
{Computed: $-2.03\times10^{-4}$ (aspect) and $-1.08\times10^{-3}$ (angle) relative, for the state real at the turn.}
{Computed in the third stage.}
{internal draft l.352, l.364; D131.}

\rChist{S9}{Photon-pair branch: ratio $1.7$ for free fields and ``about $5$'' or ``$4.96$'' with the zero-mode energy.}
{Centre $\approx5$ ($4.96$ at $\fzm=0.30$, with central amplification $d\beta_\infty/d\beta_0=26.6$) and a $90\%$ interval $3.0$--$11.9$; the earlier range $4.4$--$6.3$ was only the variation of the centre with $\fzm$, now $3.6$--$5.4$ for $\fzm\in[0.258,0.308]$ (Sec.~\ref{rep:gate}).}
{The quantum spread of the branch had not been included (T19).}
{internal draft l.42, l.256, l.598; R16 T19 (DA-5).}

\rChist{S10}{$\delta=2.04\times10^{-4}$ at $\fzm=0.30$.}
{$\delta=2.065\times10^{-4}$ (with $H_{\rm eff}$) or $2.077\times10^{-4}$, consistent with $\sigma_E=0.104$; the convention is stated where it is used.}
{Inconsistent conventions between files.}
{internal draft l.307; D125 \S5.}

\rChist{S11}{Header comment of the nonlinear survival code: pinch threshold $10^{-6}$ and ``$S_b<12h^2$''.}
{The original runs used $10^{-3}$; the criterion is stricter by the factor $0.19245$.}
{Found by the R15 check of D113.}
{C142\_HEADER\_ERRATUM.md.}

\rChist{S12}{``States that are never real survive on both branches in at most $4\%$ of cases.''}
{$4.3\%$ and $2.7\%$ at $1.2$ and $2$ times the threshold width (Wilson ranges $2.4$--$7.5\%$ and $1.3$--$5.5\%$).}
{Rounding below the computed value.}
{C153 \texttt{A2\_never\_real}; R16 consistency audit.}
\rChist{S13}{Floor values from different conventions: $0.1038$ (D109, D115), $0.1041$ (earlier draft appendix, D117, C76) and $0.1031$ (from $\delta=2.04\times10^{-4}$ in a table of the earlier draft).}
{One convention: $\sfloor=0.1037$ with $H_{\rm eff}$ and $X=2282.49$; the text writes $0.104$. This is a unification of conventions (differences below $1\%$), not a change of result.}
{Different choices of $h$ or $H_{\rm eff}$ and of $L_b$.}
{D125 \S5; C153 \texttt{B\_floor\_unification}.}
\end{description}

\subsection{The fossil in the sky}\label{sec:v23-sky}

\begin{description}
\rChist{C1}{Window and readability of the internal draft: shortest side $0.985$ to about $1.02$--$1.03\,\dLSS$; typical window $1.093$--$1.133$ ($\ell\le30$); peak readability $0.30$ near $1.07$, above $0.1$ between about $1.02$ and $1.17$; $0.45$ for free fields; $\ell\le40$ ``raises the largest readability to $0.35$''; an ensemble of nine realizations.}
{Ensemble over the plane of aspect shapes: readability is set by the shortest side; the joint probability peaks at $0.344$ at $L_0=1.075\,\dLSS$ ($0.377$ at $1.065$ at next order), exceeds $0.1$ for $L_0\in[1.018,1.183]$ and integrates to $0.042$; going from $\ell\le30$ to $\ell\le80$ raises it only by a factor $1.3$ (Sec.~\ref{rep:sky}).}
{Higher multipoles, the exact divergence, profiling over the reference volume and the full ensemble replaced the representative points.}
{internal draft l.28, 42, 174--188, 394, 404--418, App.~L; C79; D117; D125 \S4.}

\rChist{C2}{Single-fossil readout edge $1.13$--$1.14\,\dLSS$ and a factor $2.17$ in the divergence (separate line of calculation).}
{A single representative point; superseded by the ensemble (Sec.~\ref{rep:v-conditional}).}
{A ratio of divergences at one fossil is not a ratio of probabilities.}
{snapshot P5\_CMB/summary.json; D117; V12\_RESULTS\_INDEX \S D (D1$\to$D5).}

\rChist{C3}{``For $A\ge0.5$ the read edge is $1.023$--$1.031$''; ``representative interval $0.31\to0.20$, $0.45\to0.23$''; ``$P_{\rm cond}\approx1$ for $L_0\le1.03$''; ``the peak is about one third''.}
{$1.023$--$1.032$; about $0.33\to0.21$ and $0.50\to0.24$; the value near the circle limit is an extrapolation artifact (physical extrapolation $0.92$--$0.96$, lower edge $0.986$--$1.009$); the peak is $0.34$--$0.38$ depending on $\sigma_E$.}
{Independent check R15 of the ensemble layer and the wording review T12.}
{D117; R15 addendum (corrections 2, 3); T12 (W2, W3); D125 \S4.}

\rChist{C4}{Archived table of profiled divergences with a parabolic profile.}
{Pool-minimum profile over 19 reference cubes; 82 pairs had been more than $1\%$ low and one violated the data-processing inequality; the rebuilt table reproduces the peak $0.34398$ at $1.075$ (Sec.~\ref{rep:v-convergence}).}
{Found by the R15 check of the ensemble layer.}
{fossil\_table.csv $\to$ fossil\_table\_v2.csv; R15 correction 1; T12.}

\rChist{C5}{Curved cigar: ``any smooth profile''; ``the magnitude changes by less than $1\%$''; joint state: ``$\sigma_N\ge0.935$ for any $\Theta$-invariant state''; ``$L_b$ becomes larger''.}
{Profiles with constant $b$ and non-negative curvature; $1\%$ only for $R_0\ge L/2$; the bound needs a sector condition and in general reads $\sigma_N\gtrsim0.0137\,N_{\rm post}$; $L_b$ becomes smaller.}
{Independent check R15 of the separate line and wording review T12.}
{D118; D119; R15 addendum; T12 (W8, W9, W11).}

\rChist{C6}{Inflaton packet of the separate line: $\sigma_N=1.909$.}
{The width neglects the slope of the plateau; the minimal-variance value is $0.959$, with lower bound $0.935$.}
{Recomputed in our joint-state calculation.}
{D119; Sec.~\ref{rep:v-conditional}.}
\end{description}

\subsection{Failed or abandoned approaches}\label{sec:v23-failed}

\begin{description}
\item[F1] \emph{One-dimensional Born--Oppenheimer reduction of the Higgs zero mode.} Its Born--Huang term $B_w$ is almost entirely spurious near $r\to0$, and its ground energy at $\lambda=0$ is $0.151/L$ too high; replaced by the two-dimensional model (Secs.~\ref{rep:h-reduced}, \ref{rep:v-fulldim}). \emph{Files:} {\small D108; D112.}
\item[F2] \emph{Width ratios transferred between couplings.} Converting widths with ratios computed at fixed couplings was replaced by direct roots at every node; the change was only $0.001$--$0.011$~GeV, but the direct method removes the transfer error from the budget. \emph{Files:} {\small D112; D122.}
\item[F3] \emph{Tracking resonances with complex absorbing potentials alone.} In the R15 check the first eigenvalue tracking drifted to an excited resonance in 10 of 16 runs; resonances are now identified by Siegert roots with box tests, and absorbing potentials serve as an independent check. \emph{Files:} {\small R15 check of D112.}
\item[F4] \emph{Stock SMDR applications.} \texttt{calc\_fit} crashed and \texttt{calc\_all} did not finish within four hours on the research machine; a driver that calls \texttt{SMDR\_Fit\_Inputs} with the evaluation scales as arguments replaced them. \emph{Files:} {\small D120 App.~R; package \texttt{higgs/01\_sm\_mapping}.}
\item[F5] \emph{Multithreaded \texttt{mr}.} Its OpenMP build has a race condition in TSIL; it must run single-threaded. \emph{Files:} {\small R16 V1.}
\item[F6] \emph{Parabolic profile over the reference volume} (item C4) and \emph{an ensemble of nine realizations} (item C1). \emph{Files:} {\small C79; D117.}
\item[F7] \emph{Wide-step Muller iterations for the two-dimensional resonance.} Started with a wide step, the iterations converged onto Siegert zeros of box states of the $s$ direction, which gave spurious width ratios ($5971$, $1245$, $163$); the roots are now started from the interior minimum of $|f|$ on the real axis and checked with a larger box (Sec.~\ref{rep:h-reduced}). \emph{Files:} {\small D112 \S0.6, \S7.4.}
\item[F8] \emph{Sealed attempts at a finite Euclidean cap and the clock-field probes} of the separate line of calculation, which are not used (Sec.~\ref{rep:v-conditional}). \emph{Files:} {\small V12\_RESULTS\_INDEX \S B (B4, B8, B9, B13).}
\end{description}

\subsection{Statements retracted or not to be made}\label{sec:v23-retracted}

The following statements appeared in drafts, notes or discussions and are not supported by the results; the papers avoid them.
\begin{enumerate}
\item ``The Standard Model lies above the $90\%$ threshold.'' At the measured mass $\Pone(172.60)=0.84$.
\item The marginal probability, or the measured mass being on the safe side, as evidence for this beginning. The Bayes factors satisfy $|\ln B|\le\BayesMax$.
\item A probability that both branches survive. There is no decohering history probability; $P_2=\Pone^2/w_0$ is only a note (Sec.~\ref{rep:h-survival}).
\item ``Excluded by Branchina et al.'' The correct statement is conditional: incompatible with that framework.
\item The topology window ``agrees with the data'' or ``is excluded''. It is allowed by the circle and likelihood searches.
\item ``The state must be time symmetric.'' See item S2.
\item ``A stabilized Higgs''; ``new physics stabilizes''. See item H1.
\item ``The cross terms stabilize.'' See item H9.
\item ``$\pm0.17$~GeV'' as the Standard Model error; ``$s_{99}=2.2$''; ``fewer than five quanta''. See items H14, S3, S1.
\item ``A Planck-scale torus.'' The torus is a few thousand Planck lengths across, $L_b=722\,\MPbar^{-1}$.
\item ``No new physics'' without the condition on $\xiH$ and the assumption of no new states below $\mub$.
\item ``The top mass decides how the universe began''; ``tests the beginning of our universe''. The correct statement is that a measurement can exclude this beginning.
\item Specific vacuum lifetimes from Andreassen, Frost and Schwartz~\cite{Andreassen:2017rzq} (the versions of that paper quote different values), and the periodic-spinor value $-C_0$ of Guth and Vilenkin, as if they applied here.
\item ``First'', ``novel''.
\item A representative point stated as an ensemble result; a ratio of divergences stated as a ratio of probabilities.
\item Claim-layer numbers of two different model versions in one document (for example $173.41$ and $173.42$~GeV).
\item Numbers for the effect of new particles below $\mub$ on $\lambda(\mub)$ that were not computed; only assumption~\ref{assump:A2} is stated.
\end{enumerate}
